\section{Curation Protocol}

\label{sec:curation}

\subsection{Decentralized Autonomous Organization (DAO)}

Experience quality is maintained via \Zoo{} Network \DAO{}, where KEEPER token holders vote on proposed changes.

\subsubsection{Governance Token}

\textbf{KEEPER Token}:
\begin{itemize}
    \item Supply: 1,000,000,000 (fixed)
    \item Distribution: 40\% community (data contributors), 30\% foundation, 20\% team, 10\% treasury
    \item Utility: Governance voting, experience curation, compute resource allocation
\end{itemize}

\subsubsection{Voting Mechanism}

\textbf{Standard Voting}:
\begin{equation}
\text{VotingPower}(\text{voter}) = \text{KEEPER}_{\text{balance}}(\text{voter})
\end{equation}

\textbf{Quadratic Voting} (optional):
\begin{equation}
\text{VotingPower}_{\text{quad}}(\text{voter}) = \sqrt{\text{KEEPER}_{\text{balance}}(\text{voter})}
\end{equation}

Quadratic voting mitigates plutocracy by reducing marginal influence of large holders.

\subsubsection{Proposal Lifecycle}

\begin{enumerate}
    \item \textbf{Submission}: Any KEEPER holder can propose experience Add/Modify/Delete
    \item \textbf{Discussion}: 3-day community debate period
    \item \textbf{Voting}: 7-day voting window (quorum: 10\% of supply)
    \item \textbf{Execution}: If approved (66\% threshold), update registry
    \item \textbf{Archival}: Record decision on-chain and Arweave
\end{enumerate}

\subsection{Reputation System}

To incentivize high-quality contributions, we implement a reputation score:

\begin{equation}
R(\text{contributor}) = \sum_{e \in \mathcal{E}_{\text{contributor}}} w(e)
\end{equation}

where $w(e)$ is experience weight:

\begin{equation}
w(e) = e.\text{conf} \cdot \log(1 + e.\text{usage\_count}) \cdot \frac{\text{upvotes}}{\text{upvotes} + \text{downvotes}}
\end{equation}

\textbf{Benefits of High Reputation}:
\begin{itemize}
    \item Priority in proposal queue
    \item Reduced voting quorum (8\% instead of 10\%)
    \item Additional KEEPER rewards (airdropped quarterly)
\end{itemize}

\subsection{Byzantine Resistance}

\begin{definition}[Byzantine Fault Tolerance]
The Experience Ledger is $(f, n)$-Byzantine resistant if it tolerates up to $f$ malicious nodes among $n$ participants while guaranteeing:
\begin{enumerate}
    \item \textbf{Safety}: No invalid experiences enter $\mathcal{E}$
    \item \textbf{Liveness}: Valid experiences are eventually accepted
\end{enumerate}
\end{definition}

\begin{theorem}[Byzantine Robustness]
Under the assumption that at least $2f + 1$ of $n$ voters are honest, the \DAO{} voting protocol ensures safety and liveness.
\end{theorem}

\begin{proof}
\textbf{Safety}: For an invalid experience $e'$ to be accepted, it requires $> 66\%$ votes. With $2f + 1$ honest voters out of $n = 3f + 1$, at most $f$ malicious voters can collude. Thus maximum malicious votes is $f/(3f+1) < 33\%$, insufficient to approve $e'$.

\textbf{Liveness}: A valid experience $e$ requires $\geq 66\%$ votes. Honest voters ($2f+1$) represent $(2f+1)/(3f+1) \approx 67\% > 66\%$, sufficient to approve $e$ even with all malicious voters abstaining.
\end{proof}

\subsection{Sybil Resistance}

To prevent sock-puppet attacks:

\begin{enumerate}
    \item \textbf{Stake Requirement}: Minimum 1000 KEEPER to propose (economic barrier)
    \item \textbf{Identity Verification}: Optional Gitcoin Passport integration (reputation score)
    \item \textbf{Rate Limiting}: Max 3 proposals per address per week
    \item \textbf{Quadratic Voting}: Reduces impact of splitting tokens across addresses
\end{enumerate}

